% Ідентифікатори випуску очікують рішення автора; DOI прикладів не використовуються повторно.
\section{Доступність даних і коду}\label{sec:availability}

\paragraph{Дослідницькі матеріали.}
Публічний випуск для авторського огляду містить виконуваний код, збережені результати JSON, джерела рукопису, згенеровані таблиці й рисунки та обидва PDF за адресою \url{https://github.com/cyberdid/archive-index-public}, тег \code{paper-author-review-2026-10-23-r2}. DOI програмних матеріалів і препринту ще не присвоєно. Маніфест таблиць пов’язує числові матеріали з хешами джерел. Публічний пакет дає змогу перебудувати таблиці й PDF зі збережених результатів. Повний заморожений корпус, входи запитів, прогонів та оцінювання і їх 14-етапне відтворення збережені в окремому дослідницькому архіві; цей випуск їх не поширює, доки перевірка сторонньої атрибуції та сфери прав не завершена. Точне наскрізне статистичне відтворення тому потребує цих окремо збережених входів. Скрипти побудови корпусу й оцінювання оприлюднені.

\paragraph{Відтворення.}
Ізольований засіб \path{paper/run_all.sh} розрізняє складання зі збережених результатів і повне статистичне відтворення заморожених входів. Виконаний повний прохід має 14 етапів і відтворює десять звітів, канонічні qrels та записані хеші входів ANN. Фізичні вимірювання витрат, мультимодальних завдань, ANN і приватного архіву в цьому проході імпортуються, а не виконуються повторно. Зафіксовані 315 тестів реалізації, зокрема інтеграційна перевірка локально кешованої справжньої моделі, окремо характеризують поведінку коду. Разом із ним надаються описи спостереженого середовища та вимоги, зафіксовані хешами; нове повне встановлення CUDA не заявляється. Перебудова таблиць чи PDF сама по собі не повторює модельного оцінювання або фізичних вимірювань.

\paragraph{Політика джерел та атрибуція.}
Для кожного файла маніфест містить адресу джерела, зібрану ліцензійну мітку та SHA-256. Нотатки й зверстані документи походять з української та російської Вікіпедії (CC~BY-SA~4.0); 150 вибраних таблиць data.gov.ua мають мітки CC~BY. Для зображень Crossmodal-3600 збережено атрибуцію кожного зображення під CC~BY~2.0, підписи мають CC~BY~4.0; записи й транскрипти FLEURS --- CC~BY~4.0. Відбір репозиторіїв допускав мітки API MIT або Apache-2.0; зібрані репозиторії коду позначені MIT у маніфесті. Мітка репозиторію не визначає ліцензію кожного вкладеного матеріалу. Чинними залишаються сторонні повідомлення: інвентар атрибуції фіксує, зокрема, повідомлення SIL~OFL для шрифту Rubik, прив'язане до ревізії доповнення відсутньої ліцензії та питання сфери дії умов застосунку. Початкові повідомлення, файли атрибуції та аудит у \path{paper/corpus/LICENSES.md} збережено. Остаточна перевірка прав на поширення сторонніх матеріалів залишається невиконаною.

\paragraph{Ліцензування.}
Власний виконуваний код поширюється за Apache License~2.0. Текст статті, оригінальні рисунки й експериментальні записи --- за Creative Commons Attribution~4.0 International (\url{https://creativecommons.org/licenses/by/4.0/}) у межах прав автора. Посібник із ліцензування розмежовує ці матеріали та сторонні файли корпусу й залежності, для яких збережено власні умови.

\paragraph{Приватні дані та перевірка випуску.}
З архіву автора експортуються лише агреговані операційні перевірки; документи, тексти запитів, шляхи, фрагменти й індекс виключено. Перевірка випуску охоплює назви файлів, байти відомих маркерів, декодований JSON, витягнутий текст і метадані PDF та підтримуваний стиснений вміст. Чотири перевірені збіги з публічними словами/назвами потребують явних відомостей про джерело й точного хешу файла; для повних приватних запитів і домашніх шляхів винятків немає. Це не універсальна гарантія приватності й не OCR-перевірка зображень. Наявна приватна історія Git непридатна як публічний випуск; публікація потребує окремо перевіреного експорту.

\paragraph{Моделі та програмне забезпечення.}
Записи експериментів ідентифікують EmbeddingGemma~2 (\code{google/embeddinggemma-2}, ревізія \code{914f7f8}), multilingual-E5-large (\code{intfloat/multilingual-e5-large}), збірки llama.cpp/GGUF та кожну модель-суддю. Відбитки пакетів, моделей і команд залишаються прив'язаними до відповідних прогонів. Ваги моделей не входять до дослідницького експорту.

\paragraph{Участь ШІ.}
Модельне написання запитів та оцінювання релевантності описано в розділі~\ref{sec:uapar} й окремому аналізі розширення; оцінки GPT залишаються лише порівняльними. Claude і Codex допомагали реалізації та підготовці рукопису; робочий процес Scientific Agent Skills~\citep{kassis2026skills} використовувався для перевірок джерел і тверджень. Автор відповідає за план дослідження, інтерпретації, посилання та остаточні декларації. Жодна система ШІ не є автором. Ця версія не містить людського аудиту текстової релевантності; мультимодальні пари з джерельних наборів і перевірки реалізації його не замінюють.
